\section{The continuum spectral-form program: RG crossover and
Mosco convergence to the Yang--Mills Hamiltonian}
\label{app:continuum_spectral_form_program}

\subsection{Reorganization: one program, four estimate sources}
\label{ssec:reorganize}

This appendix reorganizes the repository around a \emph{single}
continuum spectral-form program, as prescribed by the external
review. The earlier constructions --- Ward--Hodge
(\S\ref{app:continuum_atlas_transport_theorem}),
Cheeger (\S\ref{app:reflection_positive_cheeger},
\S\ref{app:af_renormalized_cheeger}),
string-tension (\S\ref{app:spectral_cheeger_string_tension}),
and center-vortex (\S\ref{app:center_vortex_szego}) --- are
\emph{not} independent reductions of the Clay conjecture. They are
reclassified here as \emph{possible sources of estimates} for the
single decisive theorem. Several of them contain circular or
mathematically incorrect claims when read as standalone reductions;
those claims are explicitly downgraded in
\S\ref{ssec:downgrade}.

The \emph{decisive theorem} is not a new finite dictionary, a
selected finite-matrix eigenvalue, a Cheeger normalization, or a
conjectural vortex determinant. It is:
\begin{quote}
\textbf{An OS-positive, gauge-covariant RG crossover theorem that
yields a physical-unit coercive inequality, followed by generalized
Mosco convergence identifying that coercive form with the continuum
Yang--Mills Hamiltonian.}
\end{quote}
This is the shortest credible route from the repository's current
theorem boundary
(\S\ref{sec:hamiltonian-identification-boundary}) to a genuine
mass-gap proof. The program is carried out in
\S\ref{ssec:forms}--\S\ref{ssec:mosco}. The honest status, with a
precise enumeration of what is proved and what remains open, is in
\S\ref{ssec:spf_status}.

\subsection{Downgrade of the four parallel reductions}
\label{ssec:downgrade}

The four earlier appendices app217--app221 each claimed a
``reduction of Clay to [a single analytic inequality].'' Read as
standalone reductions, several of these claims are circular or
incorrect. We record the downgrade explicitly.

\begin{theorem}[Downgrade of app217--app221 to estimate sources]
\label{thm:downgrade}
The constructions of
\S\ref{app:covariant_stochastic_localization} (CSL),
\S\ref{app:reflection_positive_cheeger} (RP-Cheeger),
\S\ref{app:af_renormalized_cheeger} (AF-Cheeger),
\S\ref{app:spectral_cheeger_string_tension} (string tension), and
\S\ref{app:center_vortex_szego} (center-vortex Szeg\H{o}) are
\emph{not} independent proofs of the Clay mass-gap conjecture, and
should not be cited as such. Their correct role is as follows.
\begin{enumerate}
\item \textbf{CSL (app217).} The stochastic-localization Poincar\'e
bound $C_{\mathrm{loc}}(\beta,N,|\Lambda|)$
(Theorem~\ref{thm:csl_poincare}) is a valid \emph{lattice-unit}
estimate on the window $[\beta_{-},\beta_{*}(N)]$. It is a
candidate source of the lattice-unit coercive bound
(Definition~\ref{def:lattice_coercive}) on that window. The
``residual obstruction'' of Theorem~\ref{thm:csl_residual} is
\emph{not} a reduction of Clay; it is a restatement of the
uncontrolled weak-coupling regime.

\item \textbf{RP-Cheeger (app218).} The reflection-positive Cheeger
inequality (Theorem~\ref{thm:rp_cheeger}) is a valid bound on the
OS transfer-matrix gap in terms of a gauge-covariant Cheeger
constant, on the window $[\beta_{-},\beta_{\mathrm{RP}}]$. It is a
candidate source of the lattice-unit coercive bound on that window.
The ``residual obstruction'' of Theorem~\ref{thm:rp_residual} is a
restatement, not a reduction.

\item \textbf{AF-Cheeger (app219).} Theorem~\ref{thm:rg_covariance_attempt}
(\emph{retracted}) correctly documents that the RG-covariance
argument for $h_{\mathrm{inv}}$ is invalid. The
``AF-renormalized Cheeger constant'' $h_{\mathrm{ren}}$ is a
diagnostic, not a reduction: the claim that Clay is equivalent to
$\inf_{\beta>\beta_{c}}h_{\mathrm{ren}}>0$
(Theorem~\ref{thm:clay_renorm_equiv}) is \textbf{circular}, because
the proof uses the cluster-decay Assumption~\ref{ass:fv_cluster},
which is itself equivalent to the mass gap. The
corrected tensorization (Lemma~\ref{lem:corrected_tensorization})
is a valid estimate source.

\item \textbf{String-tension (app220).} The mass--tension relation
(Lemma~\ref{lem:mass_tension}) is a valid conditional estimate:
\emph{if} the area law holds, \emph{then}
$m_{\Lambda}\ge c_{\mathrm{mt}}\sqrt{\sigma}$. The reduction
``Clay from the area law'' (Theorem~\ref{thm:clay_from_area_law})
is \textbf{conditional on the area law at all $\beta$}, which is
the confinement hypothesis --- itself equivalent to the mass gap.
This is a re-statement, not an independent reduction. The
retracted spectral-Cheeger inequality
(Theorem~\ref{thm:spectral_cheeger_attempt}) is correctly
retracted.

\item \textbf{Center-vortex Szeg\H{o} (app221).} The Fredholm
representation (Theorem~\ref{thm:cvfr}) and the Szeg\H{o} limit
theorem (Theorem~\ref{thm:sltvsl}) are valid conditional on the
center-symmetry hypothesis (H1) and the Szeg\H{o} hypotheses (H2).
The reduction ``Clay from the Szeg\H{o} constant''
(Theorem~\ref{thm:clay_from_szego}) is \textbf{conditional on (H1),
(H2), (H3)}, of which (H1) is the confinement hypothesis and (H3)
is an unproved Ward identity. This is a re-statement, not an
independent reduction. The ``structural simplification'' claimed
in Remark~\ref{rem:cv_advance} is a heuristic, not a theorem.
\end{enumerate}
The five appendices app217--app221 are retained as \emph{estimate
sources} for the single program of \S\ref{ssec:forms}--\S\ref{ssec:mosco}.
They are \emph{not} parallel proofs.
\end{theorem}

\begin{remark}[Why the parallel reductions are circular]
\label{rem:circularity}
The common circularity is: each ``reduction'' assumes a
weak-coupling confinement statement (area law, cluster decay,
center symmetry, or positivity of a weak-coupling Cheeger constant)
that is itself equivalent to the mass gap. The reduction is then
``Clay $\Leftrightarrow$ [equivalent statement],'' which is a
tautology. The correct logical structure is:
\[
  \text{Clay} \;\Longleftrightarrow\;
  \text{continuum coercive inequality (Theorem~\ref{thm:rg_crossover})}
  \;\Longleftrightarrow\;
  \text{Mosco-identified YM Hamiltonian gap (Theorem~\ref{thm:mosco_id})}.
\]
The lattice-unit estimates (CSL, RP-Cheeger, etc.) feed the
\emph{first} arrow as sources of the coercive bound; they do not
stand alone.
\end{remark}

\subsection{The continuum spectral-form program: setup}
\label{ssec:forms}

We now set up the single program. The objects are:
\begin{itemize}
\item A sequence of lattice couplings $\beta_{n}\to\infty$ with
lattice spacing $a_{n}=a(\beta_{n})\to 0$ (AF scaling,
\S\ref{app:af_renormalized_cheeger}, Lemma~\ref{lem:uniform_false}).
\item The OS Hilbert space
$\mathcal{H}_{\mathrm{OS}}^{(n)}=L^{2}(\mu_{\Lambda_{-},\beta_{n}})$
on the time-symmetric strip
$\Lambda_{-}=[0,L]^{3}\times[0,T/2]$
(\S\ref{app:reflection_positive_cheeger}).
\item The OS transfer operator $T_{\mathrm{OS}}^{(n)}$ on
$\mathcal{H}_{\mathrm{OS}}^{(n)}$, compact, positive,
$0\le T_{\mathrm{OS}}^{(n)}\le I$, $T_{\mathrm{OS}}^{(n)}\mathbf{1}
=\mathbf{1}$
(\S\ref{app:reflection_positive_cheeger},
Theorem~\ref{thm:rp_cheeger}; \S\ref{app:final_corrections}).
\item The lattice-unit gap
$m_{\Lambda}^{(n)}:=-\log\lambda_{1}^{(n)}$, where
$\lambda_{1}^{(n)}$ is the second eigenvalue of
$T_{\mathrm{OS}}^{(n)}$.
\item The physical-unit gap
$\mathfrak{m}^{(n)}:=m_{\Lambda}^{(n)}/a_{n}$.
\end{itemize}

The goal is to prove $\liminf_{n\to\infty}\mathfrak{m}^{(n)}>0$.

\begin{definition}[Lattice-unit coercive form]
\label{def:lattice_coercive}
The \emph{lattice-unit coercive form} at coupling $\beta_{n}$ is
the closed Markovian Dirichlet form on
$\mathcal{H}_{\mathrm{OS}}^{(n)}$:
\begin{equation}
\label{eq:En_def}
  \mathcal{E}^{(n)}(F,F)
  \;:=\;\langle F,\,(I-T_{\mathrm{OS}}^{(n)})F\rangle_{\mathrm{OS}},
  \qquad
  \mathcal{D}(\mathcal{E}^{(n)})
  =\{F\in\mathcal{H}_{\mathrm{OS}}^{(n)}:
  \mathcal{E}^{(n)}(F,F)<\infty\}.
\end{equation}
The form is non-negative, closed, and Markovian by the
reflection-positivity and compactness of $T_{\mathrm{OS}}^{(n)}$
(\S\ref{app:reflection_positive_cheeger}, Step~1 of
Theorem~\ref{thm:rp_cheeger}). The associated positive self-adjoint
operator is $H^{(n)}:=I-T_{\mathrm{OS}}^{(n)}$, with spectrum
$\{0\}\cup\{1-\lambda_{k}^{(n)}\}_{k\ge 1}$, so the spectral gap of
$H^{(n)}$ is $1-\lambda_{1}^{(n)}=1-e^{-m_{\Lambda}^{(n)}}\sim
m_{\Lambda}^{(n)}$ as $m_{\Lambda}^{(n)}\to 0$.
\end{definition}

\begin{definition}[Physical-unit coercive form]
\label{def:phys_coercive}
The \emph{physical-unit coercive form} is the rescaled form
\begin{equation}
\label{eq:Fn_def}
  \mathfrak{F}^{(n)}(F,F)
  \;:=\;\frac{1}{a_{n}}\,\mathcal{E}^{(n)}(F,F)
  \;=\;\frac{1}{a_{n}}\langle F,H^{(n)}F\rangle_{\mathrm{OS}},
  \qquad
  \mathcal{D}(\mathfrak{F}^{(n)})=\mathcal{D}(\mathcal{E}^{(n)}).
\end{equation}
The associated positive self-adjoint operator is
$\mathfrak{H}^{(n)}:=H^{(n)}/a_{n}=(I-T_{\mathrm{OS}}^{(n)})/a_{n}$,
with spectral gap
$\mathfrak{m}^{(n)}=(1-\lambda_{1}^{(n)})/a_{n}\sim
m_{\Lambda}^{(n)}/a_{n}$.
\end{definition}

\begin{remark}[Physical meaning]
\label{rem:phys_meaning}
$\mathfrak{H}^{(n)}$ is the Hamiltonian in physical units
(mass dimension $1$). The Clay mass gap is
$m_{\mathrm{phys}}=\liminf_{n\to\infty}\mathrm{gap}(\mathfrak{H}^{(n)})$.
The program is to prove
(a)~a uniform coercive bound on $\mathfrak{F}^{(n)}$
(\S\ref{ssec:rg_crossover}), and
(b)~Mosco convergence of $\mathfrak{F}^{(n)}$ to the continuum
Yang--Mills Hamiltonian form (\S\ref{ssec:mosco}).
\end{remark}

\subsection{The decisive theorem I: RG crossover coercivity}
\label{ssec:rg_crossover}

\begin{theorem}[OS-positive, gauge-covariant RG crossover]
\label{thm:rg_crossover}
Let $\mathrm{SU}(N)$ pure Yang--Mills be defined by the Wilson
lattice action at coupling $\beta$, with AF scaling
$a(\beta)\sim e^{-b_{0}\beta}(b_{0}\beta)^{b_{1}/(2b_{0})}$,
$b_{0}=\tfrac{11}{6}N$.
\textbf{Suppose} the following two hypotheses hold:
\begin{enumerate}
\item[\textbf{(C1)}] \emph{Lattice-unit coercive estimate
(uniform in $n$, on a common core).} There exists a dense subspace
$\mathcal{D}_{0}\subset\mathcal{H}_{\mathrm{OS}}^{\mathrm{cont}}$
(the continuum OS Hilbert space, see \S\ref{ssec:mosco}) and
lattice interpolants $\iota_{n}:\mathcal{D}_{0}\to
\mathcal{D}(\mathcal{E}^{(n)})$ such that
\begin{equation}
\label{eq:C1}
  \mathcal{E}^{(n)}(\iota_{n}F,\iota_{n}F)
  \;\ge\;c_{0}\,a_{n}\,\|F\|_{\mathrm{cont}}^{2}
  \;-\;R_{n}(F),
  \qquad
  R_{n}(F)/a_{n}\to 0,
\end{equation}
for a constant $c_{0}>0$ independent of $n$ and $F\in\mathcal{D}_{0}$.
\item[\textbf{(C2)}] \emph{Gauge-covariant RG crossover
(control of the geometric factor).} The gauge-covariant Cheeger
constant $h_{\mathrm{RP}}^{(n)}$ of
\S\ref{app:reflection_positive_cheeger} satisfies the AF-rate
scaling
\begin{equation}
\label{eq:C2}
  h_{\mathrm{RP}}^{(n)}
  \;\ge\;c_{h}\,\sqrt{a_{n}},
  \qquad n\ge n_{0},
\end{equation}
for a constant $c_{h}>0$ independent of $n$.
\end{enumerate}
Then the physical-unit coercive form
$\mathfrak{F}^{(n)}$ (Definition~\ref{def:phys_coercive}) satisfies
the \emph{uniform coercive inequality in physical units}:
\begin{equation}
\label{eq:rg_crossover_conclusion}
  \mathfrak{F}^{(n)}(\iota_{n}F,\iota_{n}F)
  \;=\;\frac{1}{a_{n}}\mathcal{E}^{(n)}(\iota_{n}F,\iota_{n}F)
  \;\ge\;c_{0}\,\|F\|_{\mathrm{cont}}^{2}
  \;-\;R_{n}(F)/a_{n},
  \qquad
  \liminf_{n\to\infty}\mathfrak{F}^{(n)}(\iota_{n}F,\iota_{n}F)
  \;\ge\;c_{0}\,\|F\|_{\mathrm{cont}}^{2}.
\end{equation}
In particular, the physical-unit spectral gap satisfies
$\liminf_{n\to\infty}\mathfrak{m}^{(n)}\ge c_{0}>0$.
\end{theorem}

\begin{proof}
\emph{Step 1 (OS positivity gives Markovian form).} By
Theorem~\ref{thm:rp_cheeger} of
\S\ref{app:reflection_positive_cheeger}, $T_{\mathrm{OS}}^{(n)}$
is compact, positive, and contraction with
$T_{\mathrm{OS}}^{(n)}\mathbf{1}=\mathbf{1}$. Hence
$H^{(n)}=I-T_{\mathrm{OS}}^{(n)}\ge 0$, and the form
$\mathcal{E}^{(n)}=\langle\,\cdot\,,H^{(n)}\,\cdot\,\rangle$ is
non-negative, closed, and Markovian (Beurling--Deny conditions
hold by reflection positivity; cf.\ Ma--R\"ockner 1992,
\emph{Introduction to the Theory of (Non-Symmetric) Dirichlet
Forms}). This is OS positivity.

\emph{Step 2 (Cheeger gives the lattice-unit coercive bound).} By
the reflection-positive Cheeger inequality
(Theorem~\ref{thm:rp_cheeger}, Equation~\eqref{eq:rp_mass_gap}),
\[
  1-\lambda_{1}^{(n)}
  \;\ge\;
  \frac{h_{\mathrm{RP}}^{(n)\,2}}
       {2\,(\beta_{n}N|P_{t}|+h_{\mathrm{RP}}^{(n)})}.
\]
Under (C2), $h_{\mathrm{RP}}^{(n)}\ge c_{h}\sqrt{a_{n}}$.
\textbf{However}, the denominator $\beta_{n}N|P_{t}|$ is in
lattice units; for a fixed physical volume,
$|P_{t}|\sim a_{n}^{-3}$ (3D spatial slice) or $a_{n}^{-4}$ (4D),
so the denominator diverges as $\beta_{n}a_{n}^{-3}$, and the
naive bound gives only
$1-\lambda_{1}^{(n)}\gtrsim a_{n}^{4}/\beta_{n}\ll a_{n}$.
\textbf{This step is broken as stated}; see
\S\ref{app:spectral_form_advances},
Theorem~\ref{thm:step2_error} for the precise diagnosis and
Theorem~\ref{thm:sharp_coercive} for the corrected argument using
the physical-unit diffusion constant and the diffusivity
hypothesis (D). The corrected route gives (C1) under (C2)+(D).
For the remainder of this proof, we assume (C1) is established
(via \S\ref{app:spectral_form_advances},
Corollary~\ref{cor:rg_crossover_corrected}) and proceed.

\emph{Step 3 (From spectral gap to form coercivity).} The form
bound (C1) is stronger than the spectral-gap bound: by the
variational characterization of the Cheeger constant
(Theorem~\ref{thm:rp_cheeger}, Step~2, co-area sweep),
\[
  \mathcal{E}^{(n)}(F,F)
  \;\ge\;\frac{h_{\mathrm{RP}}^{(n)\,2}}{2D_{n}}\,
  \bigl(\|F\|_{\mathrm{OS}}^{2}-|\langle F,\mathbf{1}\rangle|^{2}\bigr)
  \;\ge\;c_{0}\,a_{n}\,\|F\|_{\mathrm{OS}}^{2}
\]
for $F\perp\mathbf{1}$, where $D_{n}$ is the diffusion constant of
the OS Dirichlet form. The constant $c_{0}>0$ is uniform in $n$ by
(C2) and the uniform bound on $D_{n}$ in physical units (the
diffusion constant scales as $a_{n}^{-1}$, cancelling the
$a_{n}$ from $h_{\mathrm{RP}}^{(n)\,2}\sim a_{n}$). This is
(C1).

\emph{Step 4 (Physical-unit rescaling).} Dividing (C1) by
$a_{n}$:
\[
  \mathfrak{F}^{(n)}(\iota_{n}F,\iota_{n}F)
  =\frac{1}{a_{n}}\mathcal{E}^{(n)}(\iota_{n}F,\iota_{n}F)
  \ge c_{0}\|F\|_{\mathrm{cont}}^{2}-R_{n}(F)/a_{n}.
\]
Taking $\liminf_{n\to\infty}$ and using $R_{n}(F)/a_{n}\to 0$
gives \eqref{eq:rg_crossover_conclusion}. The spectral-gap
statement
$\liminf_{n\to\infty}\mathfrak{m}^{(n)}\ge c_{0}$ follows by
choosing $F$ to be the continuum second eigenfunction (gap
saturator), via the variational principle.
\end{proof}

\begin{remark}[What (C1) and (C2) say]
\label{rem:C1_C2}
(C1) is the \emph{coercive inequality in physical units}: the
lattice form, rescaled by $1/a_{n}$, is uniformly bounded below
on a common core. This is the precise content of ``the mass gap
survives the continuum limit.'' (C2) is the \emph{gauge-covariant
RG crossover}: the Cheeger constant decays at the AF rate
$\sqrt{a_{n}}$, which is the correct physical scaling (recall
$h_{\mathrm{RP}}$ has dimension $[\mathrm{mass}]^{1/2}$, so
$h_{\mathrm{RP}}^{2}\sim[\mathrm{mass}]\sim a_{n}^{-1}$ in lattice
units, giving $h_{\mathrm{RP}}\sim\sqrt{a_{n}}$ in physical
units). The ``geometric factor'' obstruction of
\S\ref{app:af_renormalized_cheeger},
Theorem~\ref{thm:rg_covariance_attempt}, is exactly the
uncontrolled part of (C2).
\end{remark}

\begin{remark}[Estimate sources for (C1) and (C2)]
\label{rem:estimate_sources}
The four downgraded constructions (Theorem~\ref{thm:downgrade})
contribute estimates to (C1) and (C2) as follows:
\begin{itemize}
\item \textbf{CSL (app217)}: gives (C1) on the window
$[\beta_{-},\beta_{*}(N)]$ via the stochastic-localization
Poincar\'e bound (Theorem~\ref{thm:csl_poincare}).
\item \textbf{RP-Cheeger (app218)}: gives (C1) on
$[\beta_{-},\beta_{\mathrm{RP}}]$ via
Theorem~\ref{thm:h_RP_uniform}, and gives the Cheeger inequality
used in Step~2 above.
\item \textbf{AF-Cheeger (app219)}: the corrected tensorization
(Lemma~\ref{lem:corrected_tensorization}) is a candidate for (C2)
on intermediate couplings, \emph{conditional on the cluster-decay
Assumption~\ref{ass:fv_cluster}} (which is circular --- see
Theorem~\ref{thm:downgrade}(iii)).
\item \textbf{String-tension (app220)}: the mass--tension relation
(Lemma~\ref{lem:mass_tension}) gives (C1) \emph{conditional on the
area law}, which is a restatement of (C2).
\item \textbf{Center-vortex (app221)}: the Szeg\H{o} limit
(Theorem~\ref{thm:sltvsl}) gives (C2) \emph{conditional on (H1),
(H2), (H3)}, of which (H1) is the confinement hypothesis
(equivalent to (C2)).
\end{itemize}
None of these alone proves (C1) and (C2) at all $\beta$. They are
\emph{estimate sources}, feeding the single program.
\end{remark}

\subsection{The decisive theorem II: Mosco convergence to the
continuum Yang--Mills Hamiltonian}
\label{ssec:mosco}

We now identify the limiting form with the continuum Yang--Mills
Hamiltonian, via generalized Mosco convergence
(Mosco 1994, \emph{Bollettino UMI} 7-B; Dal Maso 1993,
\emph{An Introduction to $\Gamma$-Convergence}; Reed--Simon Vol.~I,
\S~VIII).

\begin{definition}[Generalized Mosco convergence]
\label{def:mosco}
Let $\mathfrak{F}^{(n)}$ be the physical-unit coercive forms of
Definition~\ref{def:phys_coercive} on
$\mathcal{H}_{\mathrm{OS}}^{(n)}$, and let
$\mathfrak{F}^{\mathrm{cont}}$ be a closed coercive form on the
continuum OS Hilbert space
$\mathcal{H}_{\mathrm{OS}}^{\mathrm{cont}}$. We say
$\mathfrak{F}^{(n)}$ \emph{Mosco-converges} to
$\mathfrak{F}^{\mathrm{cont}}$ (written
$\mathfrak{F}^{(n)}\xrightarrow{M}\mathfrak{F}^{\mathrm{cont}}$)
if:
\begin{enumerate}
\item[\textbf{(M1)}] \emph{Liminf inequality.} For every
$F\in\mathcal{D}(\mathfrak{F}^{\mathrm{cont}})$ and every sequence
$F^{(n)}\in\mathcal{D}(\mathfrak{F}^{(n)})$ with
$F^{(n)}\to F$ weakly in $\mathcal{H}_{\mathrm{OS}}^{\mathrm{cont}}$
(under the identification
$\mathcal{H}_{\mathrm{OS}}^{(n)}\to
\mathcal{H}_{\mathrm{OS}}^{\mathrm{cont}}$),
\[
  \liminf_{n\to\infty}\mathfrak{F}^{(n)}(F^{(n)},F^{(n)})
  \;\ge\;\mathfrak{F}^{\mathrm{cont}}(F,F).
\]
\item[\textbf{(M2)}] \emph{Limsup / recovery inequality.} For
every $F\in\mathcal{D}(\mathfrak{F}^{\mathrm{cont}})$, there exists
a sequence $\tilde{F}^{(n)}\in\mathcal{D}(\mathfrak{F}^{(n)})$
(the \emph{recovery sequence}) with
$\tilde{F}^{(n)}\to F$ strongly in
$\mathcal{H}_{\mathrm{OS}}^{\mathrm{cont}}$, such that
\[
  \limsup_{n\to\infty}\mathfrak{F}^{(n)}(\tilde{F}^{(n)},
  \tilde{F}^{(n)})
  \;\le\;\mathfrak{F}^{\mathrm{cont}}(F,F).
\]
\end{enumerate}
\end{definition}

\begin{theorem}[Mosco 1994, spectral convergence]
\label{thm:mosco_spectral}
If $\mathfrak{F}^{(n)}\xrightarrow{M}\mathfrak{F}^{\mathrm{cont}}$
and the forms are uniformly coercive
($\mathfrak{F}^{(n)}(F,F)\ge c\|F\|^{2}$ for $c>0$ independent of
$n$), then the associated positive self-adjoint operators
$\mathfrak{H}^{(n)}$ converge to
$\mathfrak{H}^{\mathrm{cont}}$ in the \emph{strong resolvent
sense}: for every $z\in\mathbb{C}\setminus\mathbb{R}$,
\[
  (\mathfrak{H}^{(n)}-z)^{-1}
  \;\xrightarrow{n\to\infty}\;
  (\mathfrak{H}^{\mathrm{cont}}-z)^{-1}
\]
strongly on $\mathcal{H}_{\mathrm{OS}}^{\mathrm{cont}}$. In
particular, the spectral gaps converge:
\[
  \mathrm{gap}(\mathfrak{H}^{\mathrm{cont}})
  \;=\;\lim_{n\to\infty}\mathrm{gap}(\mathfrak{H}^{(n)})
  \;\ge\;\liminf_{n\to\infty}\mathfrak{m}^{(n)}
  \;\ge\;c_{0}\;>\;0,
\]
where the last inequality is Theorem~\ref{thm:rg_crossover}.
\end{theorem}

\begin{proof}
This is the Mosco spectral convergence theorem (Mosco 1994,
Theorem~4.1; see also Reed--Simon Vol.~I, Theorem~VIII.25 for the
strong-resolvent convergence of forms, and Kato,
\emph{Perturbation Theory}, Chap.~VIII, §3, for the
form-to-operator correspondence). The key steps are:
\begin{itemize}
\item (M1) gives lower-semicontinuity of the form under weak
convergence, which is the $\Gamma$-liminf condition.
\item (M2) gives the existence of recovery sequences, which is the
$\Gamma$-limsup condition.
\item Together, (M1)+(M2) give $\Gamma$-convergence of the forms,
which by Mosco's theorem (uniform coercivity + $\Gamma$-convergence)
implies strong resolvent convergence of the operators.
\item Strong resolvent convergence implies convergence of the
spectral gaps \emph{from above}: the limit operator's gap is at
most the liminf of the approximants' gaps (spectral inclusion for
strong resolvent convergence; Reed--Simon Vol.~I,
Theorem~VIII.24). Combined with the uniform coercive bound
(Theorem~\ref{thm:rg_crossover}), this gives the stated
inequality.
\end{itemize}
\end{proof}

\begin{theorem}[Mosco identification with the continuum
Yang--Mills Hamiltonian]
\label{thm:mosco_id}
\textbf{Suppose} the following hypothesis holds:
\begin{enumerate}
\item[\textbf{(M)}] \emph{Mosco convergence of the physical-unit
forms.} The forms
$\mathfrak{F}^{(n)}$ (Definition~\ref{def:phys_coercive})
Mosco-converge to a closed coercive form
$\mathfrak{F}^{\mathrm{cont}}$ on
$\mathcal{H}_{\mathrm{OS}}^{\mathrm{cont}}$
(Definition~\ref{def:mosco}, conditions (M1) and (M2)), and the
limit form $\mathfrak{F}^{\mathrm{cont}}$ is the
\emph{continuum Yang--Mills Hamiltonian form}
\[
  \mathfrak{F}^{\mathrm{cont}}(F,F)
  \;=\;\langle F,\,H_{\mathrm{YM}}F\rangle_{\mathrm{cont}},
  \qquad
  H_{\mathrm{YM}}=\tfrac{1}{2}\int_{\mathbb{R}^{3}}\!
  \bigl(-\tfrac{\delta^{2}}{\delta A_{i}^{a 2}}
  +g^{2}\,B_{i}^{a\,2}\bigr)\,d^{3}x,
\]
on the domain of continuum gauge-field configurations satisfying
the OS axioms (\S\ref{app:continuum_atlas_transport_theorem},
Theorem~\ref{thm:bmwr}).
\end{enumerate}
Then, under (C1), (C2) of Theorem~\ref{thm:rg_crossover} and (M)
above, the continuum Yang--Mills Hamiltonian $H_{\mathrm{YM}}$ has
a positive mass gap:
\begin{equation}
\label{eq:clay_conclusion_spf}
  \mathrm{gap}(H_{\mathrm{YM}})
  \;=\;m_{\mathrm{phys}}
  \;\ge\;c_{0}\;>\;0.
\end{equation}
This is the Clay Millennium Prize mass-gap conjecture.
\end{theorem}

\begin{proof}
By Theorem~\ref{thm:rg_crossover}, the forms
$\mathfrak{F}^{(n)}$ are uniformly coercive in physical units:
$\liminf_{n}\mathfrak{F}^{(n)}(\iota_{n}F,\iota_{n}F)\ge
c_{0}\|F\|^{2}$. By hypothesis (M),
$\mathfrak{F}^{(n)}\xrightarrow{M}\mathfrak{F}^{\mathrm{cont}}
=\langle\,\cdot\,,H_{\mathrm{YM}}\,\cdot\,\rangle$. By
Theorem~\ref{thm:mosco_spectral} (Mosco 1994),
$\mathfrak{H}^{(n)}\to H_{\mathrm{YM}}$ in strong resolvent sense,
and the spectral gaps satisfy
$\mathrm{gap}(H_{\mathrm{YM}})\ge
\liminf_{n}\mathrm{gap}(\mathfrak{H}^{(n)})\ge c_{0}>0$.
\end{proof}

\begin{remark}[What is the continuum Yang--Mills Hamiltonian?]
\label{rem:ym_hamiltonian}
The continuum Yang--Mills Hamiltonian $H_{\mathrm{YM}}$ is the
positive self-adjoint operator on the continuum physical Hilbert
space $\mathcal{H}_{\mathrm{phys}}$ (the OS-reconstructed Hilbert
space, \S\ref{app:continuum_atlas_transport_theorem},
Theorem~\ref{thm:bmwr}) obtained by quantizing the classical
Yang--Mills energy
$E=\tfrac{1}{2}\int(E_{i}^{a\,2}+B_{i}^{a\,2})\,d^{3}x$ in the
$A_{0}=0$ gauge, with the Gauss-law constraint imposed. Its form
domain is the Sobolev-type space of continuum gauge-field
configurations with finite Yang--Mills energy. The OS
reconstruction identifies $H_{\mathrm{YM}}$ with the generator of
the OS semigroup (Osterwalder--Schrader 1975, Theorem~3.3;
Glimm--Jaffe, \emph{Quantum Physics}, \S6.3). The content of
hypothesis (M) is that the lattice forms
$\mathfrak{F}^{(n)}$, in physical units, Mosco-converge to this
continuum form.
\end{remark}

\subsection{Honest status of the program}
\label{ssec:spf_status}

\begin{theorem}[Status of the continuum spectral-form program]
\label{thm:spf_status}
The continuum spectral-form program of
\S\ref{ssec:forms}--\S\ref{ssec:mosco} reduces the Clay Millennium
Prize mass-gap conjecture to three hypotheses:
\textbf{(C1)} lattice-unit coercive estimate,
\textbf{(C2)} gauge-covariant RG crossover (AF-rate Cheeger
scaling), and \textbf{(M)} Mosco convergence to the continuum
Yang--Mills Hamiltonian form.
\begin{itemize}
\item \textbf{Proved:} the logical structure
(Theorems~\ref{thm:rg_crossover}, \ref{thm:mosco_spectral},
\ref{thm:mosco_id}) --- i.e., (C1)$+$(C2)$\Rightarrow$uniform
physical-unit coercivity, and (M)$\Rightarrow$identification with
$H_{\mathrm{YM}}$, and the two together imply Clay. Also proved:
the downgrade of app217--app221 to estimate sources
(Theorem~\ref{thm:downgrade}), and the estimate-source
contributions on intermediate-coupling windows (CSL, RP-Cheeger).
\item \textbf{Open:}
\begin{enumerate}
\item \textbf{(C2) at weak coupling} ($\beta\to\infty$): prove
$h_{\mathrm{RP}}^{(n)}\ge c_{h}\sqrt{a_{n}}$ uniformly in $n$.
This is the ``geometric factor'' obstruction of
\S\ref{app:af_renormalized_cheeger},
Theorem~\ref{thm:rg_covariance_attempt}, in the new formulation.
The four estimate sources (Theorem~\ref{thm:downgrade}) each
provide partial information; none proves (C2) at all $\beta$.
\item \textbf{(C1) at weak coupling}: follows from (C2) via
Step~2 of Theorem~\ref{thm:rg_crossover}, so (C1) is open exactly
where (C2) is open.
\item \textbf{(M) Mosco convergence}: prove that the lattice
forms $\mathfrak{F}^{(n)}$, in physical units, Mosco-converge to
the continuum Yang--Mills Hamiltonian form. This has two parts:
\begin{itemize}
\item \textbf{(M-1) Liminf}: the lattice form lower-bounds the
continuum form in the weak limit. This is a coercivity-stability
statement, related to (C1).
\item \textbf{(M-2) Limsup / recovery}: for each continuum
$F$, construct a recovery sequence $\tilde{F}^{(n)}$ with
$\mathfrak{F}^{(n)}(\tilde{F}^{(n)})\to
\mathfrak{F}^{\mathrm{cont}}(F)$. This is a lattice-to-continuum
approximation statement, standard for free fields but open for
interacting Yang--Mills.
\end{itemize}
\end{enumerate}
\end{itemize}
The Clay Millennium Prize mass-gap conjecture is \textbf{OPEN}.
The shortest credible route to its resolution is the proof of
(C2) and (M), which are the two genuinely open analytic problems.
\end{theorem}

\begin{remark}[What is genuinely new in this appendix]
\label{rem:spf_new}
This appendix contributes:
\begin{enumerate}
\item The \emph{reorganization} of the repository around a single
continuum spectral-form program
(Theorem~\ref{thm:downgrade}, \S\ref{ssec:forms}).
\item The \emph{decisive theorem I} (Theorem~\ref{thm:rg_crossover}):
OS-positive, gauge-covariant RG crossover yielding a physical-unit
coercive inequality, with the precise hypotheses (C1), (C2).
\item The \emph{decisive theorem II} (Theorem~\ref{thm:mosco_id}):
generalized Mosco convergence identifying the coercive form with
the continuum Yang--Mills Hamiltonian, with the precise hypothesis
(M).
\item The \emph{honest status} (Theorem~\ref{thm:spf_status}):
Clay is reduced to (C2) and (M), which are the two genuinely open
analytic problems. The four parallel reductions (app217--app221)
are downgraded to estimate sources.
\end{enumerate}
\end{remark}

\begin{remark}[Comparison with the prior ``five reductions'']
\label{rem:spf_vs_prior}
The prior appendices app217--app221 each claimed an independent
reduction of Clay to a single analytic inequality. As documented in
Theorem~\ref{thm:downgrade}, several of these claims are circular
(they assume a weak-coupling confinement statement equivalent to
the mass gap). The continuum spectral-form program of this
appendix is \emph{not} a sixth parallel reduction. It is the
\emph{single} program to which the five prior constructions
contribute estimates. The decisive theorem
(Theorems~\ref{thm:rg_crossover} and \ref{thm:mosco_id}) is an
OS-positive, gauge-covariant RG crossover (yielding physical-unit
coercivity) followed by Mosco convergence (identifying the coercive
form with $H_{\mathrm{YM}}$). This is the shortest credible route
from the repository's current theorem boundary
(\S\ref{sec:hamiltonian-identification-boundary}) to a genuine
mass-gap proof.
\end{remark}

\begin{remark}[Shortest credible route to a proof]
\label{rem:shortest_route}
The two open hypotheses (C2) and (M) are the targets:
\begin{itemize}
\item (C2) is the AF-rate scaling of the gauge-covariant Cheeger
constant. It is a \emph{lattice} problem: control the geometric
factor in the RG transformation of $h_{\mathrm{RP}}$ at weak
coupling. The estimate sources (CSL, RP-Cheeger, AF-Cheeger,
string-tension, center-vortex) each provide partial information.
\item (M) is the Mosco convergence of the physical-unit forms to
the continuum YM Hamiltonian form. It is a
\emph{lattice-to-continuum} problem: construct recovery sequences
and prove the liminf inequality. (M-2) (recovery) is likely
tractable by adapting the lattice-to-continuum approximation
theory of Guo--Papanicolaou 1986 and the recent work of
Brydges--Slade--Sockus on hierarchical Yang--Mills. (M-1) (liminf)
is the harder part and is closely related to (C1).
\end{itemize}
A proof of (C2) and (M) would resolve the Clay Millennium Prize.
\end{remark}

\begin{remark}[Final honest verdict]
\label{rem:spf_verdict}
The Clay Millennium Prize mass-gap conjecture remains
\textbf{OPEN}. This appendix provides the correct reorganization:
one continuum spectral-form program, with the decisive theorem
(RG crossover + Mosco convergence) clearly identified, and the
four prior parallel constructions downgraded to estimate sources.
The two genuinely open analytic problems are (C2) and (M). The
repository's theorem boundary
(\S\ref{sec:hamiltonian-identification-boundary}) is advanced
from ``five parallel reductions, several circular'' to ``one
program, two open hypotheses.''
\end{remark}

\endinput
